\documentclass[12pt]{article}

\usepackage[utf8]{inputenc}
\usepackage[spanish, es-noshorthands]{babel}

\usepackage{mathtools}
\usepackage{amssymb}
\usepackage{amsthm}

\usepackage{geometry}
\geometry{margin=2.5cm}

\usepackage{tikz}
\usetikzlibrary{positioning}

\usepackage{hyperref}

\title{Entregables Teoría de Números}
\author{Juan Felipe Rodríguez Córdoba}
\date{}

\begin{document}
\maketitle

\section*{Ejercicio 1}
Se sabe que
\[
3 \mid 4^n - 1.
\]
Demuestra que
\[
9 \mid 4^n - 3n - 1
\]

\vspace{8cm}

\section*{Ejercicio 2}
¿Es verdad que para cualquier \(n \in \mathbb{N}\) resulta que
\[
P(n)=n^2+n+41
\]
es primo?

\vspace{5cm}

\section*{Ejercicio 3}
¿Cuántos divisores tienen estos números?
\[
11\cdot 13,
\]
\[
11\cdot 13\cdot 17,
\]
\[
11^{14},
\]
\[
11^{17}\cdot 13^{28}\cdot 17^{65}.
\]

\vspace{8cm}

\section*{Ejercicio 4}
Demuestra que para cualquier primo \(p>3\) se cumple que
\[
p^2-1
\]
es divisible por \(24\).

\vspace{5cm}

\section*{Ejercicio 5}
Demuestra que
\[
a^3+b^3+c^3-3abc
\]
es divisible por
\[
a+b+c.
\]

\vspace{7cm}

\section*{Ejercicio 6}
Demuestra que para cualquier \(n\in\mathbb{N}\) se cumple que
\[
k! \mid n(n+1)(n+2)\cdots(n+k-1).
\]

\vspace{7cm}

\section*{Ejercicio 13}
Demuestra que si
\[
\gcd(m,n)=1,
\]
entonces
\[
a\equiv b \pmod{mn}
\]
si y solo si
\[
a\equiv b \pmod m
\]
y
\[
a\equiv b \pmod n.
\]

\vspace{7cm}

\section*{Ejercicio 16}
En una tabla cuadrada hemos coloreado una tabla cuadrada más pequeña.
Han quedado \(79\) casillas sin colorear.
¿Pueden quedar las esquinas sin colorear?

\vspace{7cm}

\section*{Ejercicio 1}
Representa
\[
\sqrt{7}
\]
como fracción continua periódica.

\vspace{20cm}

\section*{Ejercicio 2}
Demuestra que hay infinitos primos del tipo
\[
4n+3.
\]

\vspace{20cm}

\section*{Ejercicio 3}
Busca todos los números de \(3\) cifras cuyo cuadrado termina en este propio número.

\vspace{20cm}

\section*{Ejercicio 4}
Demuestra que la ecuación
\[
4x^2+3y^2=2000000002
\]
no tiene solución en números enteros.

\vspace{20cm}

\section*{Ejercicio 7}
Dado el mensaje cifrado
\[
\texttt{t5r](+},
\]
descifra el mensaje si para cifrarlo se utilizó el algoritmo RSA con
\[
n=77,\qquad e=37,
\]
y la equivalencia alfanumérica dada en la hoja.

\end{document}